\section*{A route through the chapter}
\addcontentsline{toc}{section}{A route through the chapter}
The central distinction is between a residual root, a coherent branch and an
occupied packet in a carrier. Sections 2--7 develop the finite and formal
machinery. Sections 8--11 assemble its arithmetic uses under stated hypotheses.
The residue-counting arguments appear in Appendix C. Appendix D states the
conditional occupied-system completion theorem. The visual studies provide
worked examples of the same constructions.

\begin{tabularx}{\linewidth}{@{}p{.24\linewidth}X@{}}\toprule
Read for & Follow this route\\\midrule
The mathematical core & Residual partition, reconstruction, smooth and singular counts, cohomology packets and directed seed rank.\\
The arithmetic assembly & Common coefficient data, finite Frobenius separation, explicit carrier summands and inverse-limit-closed profiles.\\
The formal contract & The formal companion explains the terminal theorem and its three direct dependency modules, with a statement register and a hypothesis-by-hypothesis guide.\\
The finite examples & Plates X.1--X.30 show residue tables, projector tests, countermodels and norm-character examples.\\\bottomrule
\end{tabularx}

\paragraph{Two outputs require two sets of inputs.}
The terminal theorem proves that a transported target family exists and that the
declared summand can be recovered inside a supplied compatible carrier. It does
not identify those two objects. An application that needs such an identification
must supply its comparison. Likewise, zero multiplicities are permitted by the
formal assembly, so a nonzero conclusion retains a separate nonzero premise.

\paragraph{Mathematical and formal scope.}
The final source is over the rational absolute Galois group. Carrier data,
constituent comparisons and a completed source family are explicit inputs.
The formal companion distinguishes the amended conditional statements from
stronger existence questions and identifies their recorded verification.
Earlier source contracts remain available in the accompanying research records.
